\documentclass[10pt,a4paper]{amsart}
\usepackage[margin=19mm]{geometry}
\usepackage{amsmath,amssymb,mathtools}
\usepackage[scr=rsfs,cal=euler]{mathalfa}
\usepackage{xfrac}
\usepackage[unicode,hidelinks,bookmarks=false]{hyperref}
\newcommand{\ie}{i.e.\ }
\newcommand{\cf}{cf.\ }
\newcommand{\resp}{respectively\ }
\newcommand{\Subsection}{\subsection}
\providecommand{\nobreakdash}{\nobreak\hbox{-}}
\setlength{\emergencystretch}{2em}
\multlinegap0pt
\usepackage{bm}
\usepackage{bbm}
\usepackage{dsfont}
\usepackage{xspace}
\usepackage{cancel}
\usepackage{mathtools}
\usepackage{amsthm}
\makeatletter
\@ifundefined{theorem}{\newtheorem{theorem}{Theorem}}{}
\@ifundefined{prop}{\newtheorem{prop}[theorem]{Proposition}}{}
\makeatother
\newcommand{\CP}{\mathbb{CP}}
\newcommand{\C}{\mathbb{C}}
\newcommand{\cX}{\mathcal{X}}
\title[Maximal twisted Betti numbers of complex hyperplane arrangem]{Maximal twisted Betti numbers of complex hyperplane arrangement complements}
\author{Yongqiang Liu, Laurentiu Maxim, Botong Wang}
\date{Source: arXiv:2511.08787v1 (CC BY 4.0)}
\begin{document}
\maketitle

\noindent\emph{Source and licence.} Maximal twisted Betti numbers of complex hyperplane arrangement complements. Version arXiv:2511.08787v1 (CC BY 4.0), distributed under the Creative Commons Attribution 4.0 International licence, as recorded in the arXiv licence metadata for this exact version and shown on the abstract page https://arxiv.org/abs/2511.08787v1. Full rights notice: licenses/arxiv-2511.08787-CC-BY-4.0.txt.\par\smallskip

\noindent\textit{Editorial scope.} Counted unit: the Prop whose source label is \texttt{propA} in the article (source0.tex lines 415--421, source label \texttt{propA}), delivered together with the article's own complete original proof of it (source0.tex lines 422--438, one \texttt{proof} environment of 17 lines). No mathematics is written, repaired or re-derived: statement and proof are the article's own text line for line. Required context retained uncounted: none. Every definition, display and label of those lines is retained unchanged. Omitted from this package and not redistributed: 1--414, 439--605. Nothing inside a retained excerpt is omitted or altered: every retained body is delivered exactly as the article's lines are written, so no omission receipt of any kind accompanies this package. Citations: the retained excerpts carry no citation command at all, so no bibliography entry of the article is transcribed and no reference is redistributed. Reference adaptation: none was needed and none was made; every reference inside the delivered unit resolves inside it, so no unresolved reference or citation is produced. Printed slips kept literal: no printed word, symbol, hypothesis or number of the counted statement or of its proof was altered. Layout and class: the article's own class is \texttt{amsart} with packages including geometry, ulem, amsmath, amscd, amssymb, amsfonts, latexsym, wasysym, mathrsfs, mathtools, hhline, color, tikz, xy; the wrapper is the lane's pinned \texttt{amsart} editorial wrapper at 10pt on A4, which restates the theorem-like environments the retained text needs and adds the article's own macro definitions that the retained text needs (\texttt{\textbackslash C}, \texttt{\textbackslash CP}, \texttt{\textbackslash cX}), with extra wrapper packages (bm, bbm, dsfont, xspace, cancel, mathtools). The delivered numbering is the unit's own. Assets: no retained span contains a figure environment, a graphics-inclusion command, a PDF-page inclusion, a table or a TikZ picture; no image file is copied into this package, the package declares no asset dependency and no image file is redistributed with it. Licence: version arXiv:2511.08787v1 of this article is distributed under the Creative Commons Attribution 4.0 International licence, as recorded in the arXiv licence metadata for this exact version and shown on the abstract page https://arxiv.org/abs/2511.08787v1. A full rights notice is delivered at licenses/arxiv-2511.08787-CC-BY-4.0.txt. Characters: every retained excerpt is pure ASCII, so the delivered engine, XeLaTeX with its default Latin Modern text font, reports no missing character.
\begin{prop}\label{propA}
Under the above notations, for a general linear form $l$, the following statements hold.
\begin{enumerate}
\item The pull-back of the Whitney stratification $\mathcal{W}$ of $\CP^n$ by $p$ defines a Whitney stratification of $X$, which we denote by $\mathcal{W}'$. 
\item The map $\pi_X\colon X\to \C$ does not have critical points along $E$, with respect to the stratification $\mathcal{W}'$. 
\end{enumerate}
\end{prop}

\begin{proof}
Notice that $[x_0:l]$ defines a rational map $\CP^n\dashrightarrow \CP^1$ and $\cX$ is equal to the closure of its graph in $\CP^n\times \CP^1$. Moreover, $X$ is equal to the intersection of $\cX$ with $\CP^n\times \C$. 

Taking cartesian product with $\C$, the Whitney stratification $\mathcal{W}$ of $\CP^n$ induces a Whitney stratification $\mathcal{W}_{\CP^n\times \C}$ on $\CP^n\times \C$. For the first statement, we will show that $X$ intersects every stratum of $\mathcal{W}_{\CP^n\times \C}$ transversally, which implies that the restriction of $\mathcal{W}_{\CP^n\times \C}$ to $X$ is a Whitney stratification. 

Let $t$ be the coordinate of $\C$. Then $X \subset \CP^n\times \C$ is defined by the equation $l=t \cdot x_0$, or
\[
a_0x_0+a_1x_1+\cdots +a_nx_n-t x_0=0
\]
for general $a_0, \dots, a_n$. Let $\mathcal{O}_{\CP^n\times \C}(1)$ be the pullback of the line bundle $\mathcal{O}_{\CP^n}(1)$ by the projection. Then, via pullback, $x_0, \dots, x_n$ define global sections of $\mathcal{O}_{\CP^n\times \C}(1)$. Also via pullback, $t$ is a global function of $\CP^n\times \C$, and hence the product $t x_0$ is also a global section of $\mathcal{O}_{\CP^n\times \C}(1)$. Notice that the global sections $x_0, \dots, x_n, tx_0$ define a base-point free linear system of $\mathcal{O}_{\CP^n\times \C}(1)$, and $X$ is  a general member of this linear system. Therefore, by the Bertini theorem for basepoint-free linear systems, the zero locus of a general section is transversal to the Whitney stratification $\mathcal{W}_{\CP^n\times \C}$. Hence the first statement follows. 

For the second statement, we need to show that for any $x\in \C$, $(\pi_X)^{-1}(x)$ intersects every stratum of $\mathcal{W}'$ transversally along $E$. Notice that 
\[
E= V_l\times \C\subset \CP^n\times \C.
\]
Let $\tilde{x}$ be any point in $(\pi_X)^{-1}(x)\cap E$. Then $\tilde{x}$ is contained in the line $p(\tilde{x})\times \C$ in $E$. Since $\mathcal{W}'$ is defined as the pullback of a stratification of $\CP^n$, the line $p(\tilde{x})\times \C$ is contained in one stratum of $\mathcal{W}'$. Since the line $p(\tilde{x})\times \C$ intersects $(\pi_X)^{-1}(x)$ transversally, the stratum of $\mathcal{W}'$ containing $\tilde{x}$ also intersects $(\pi_X)^{-1}(x)$ transversally. Therefore, the intersection of  $(\pi_X)^{-1}(x)$ and any stratum of $\mathcal{W}'$ is transversal along $E$. Thus, the second statement follows. 
\end{proof}

\end{document}
